\chapter{SVD Sigma-Point Filters}
\label{ch:bf-svd-sigma-point}

The SVD sigma-point filter is the most important nonlinear-filtering lesson
from the recent source-project debugging work.  It improved value-side
robustness by using spectral covariance factors and explicit finite-arithmetic
guards, but it also exposed why raw autodiff through spectral decompositions is
not a comfortable industrial HMC default.  BayesFilter should preserve both
lessons.

\section{What the Source Audit Established}
\label{sec:bf-svd-audit-established}

The math/code audit recorded in
the source-project SVD filter audit result from 2026-04-25
established the following bounded facts:
\begin{itemize}
  \item the linear SVD Kalman path uses a Joseph-style covariance update;
  \item the nonlinear sigma-point path uses the simplified update
    \[
      P^+ = P^- - K P_{zz} K^\top ;
    \]
  \item PSD projection occurs after the candidate covariance is formed, so it
    cannot prevent overflow or non-finite values during the formation of
    $P^+$;
  \item the conditioning guard with threshold
    \texttt{\_SVD\_CONDITIONING\_MAX\_ABS = 1e100} is an engineering
    finite-arithmetic guard, not a mathematically derived posterior threshold;
  \item prior short HMC diagnostics were numerical-path regressions, not
    convergence evidence.
\end{itemize}
These are implementation facts, not general theorems about SVD filtering.
BayesFilter should cite them as project evidence and should avoid converting
them into unsupported production-readiness claims.

\section{Spectral Differentiation Risk}
\label{sec:bf-svd-spectral-risk}

The Phase 2B literature gate accepted the Matrix Backprop source for the
bounded claim that gradients through SVD or eigendecomposition contain terms
with denominators involving singular-value or eigenvalue gaps.  When two
singular values or eigenvalues are equal, the decomposition is not uniquely
differentiable; when they are merely close, derivative terms can become large
and numerically unstable.

This is directly relevant to HMC.  A large DSGE or NAWM-sized model can visit
parameter regions where covariance factors are ill conditioned, nearly rank
deficient, or spectrally clustered.  A value-side SVD factor may still return a
finite covariance approximation, while the tape gradient through the spectral
basis becomes noisy, discontinuous, or non-finite.  Larger dimensions make
close spectral gaps more likely in practice, so BayesFilter should require
gap telemetry and derivative validation before using SVD sigma-point gradients
inside production HMC.

\section{Analytic Score and Hessian Recursion}
\label{sec:bf-svd-sp-analytic-recursion}

The SVD sigma-point derivative problem has two layers.  The likelihood layer is
not new: after the sigma-point update has produced
\[
  v_t,\quad S_t,\quad
  \dot v_t^{(i)},\quad \dot S_t^{(i)},\quad
  \ddot v_t^{(ij)},\quad \ddot S_t^{(ij)},
\]
the score and Hessian are exactly the solve-form expressions in
Chapter~\ref{ch:bf-kalman-hessian}.  The new work is the sigma-point layer:
differentiate the point cloud, the nonlinear maps, and the weighted moment
assemblies that generate these innovation objects.

Thus this chapter plugs into the analytic-derivative stack as follows.
Chapter~\ref{ch:bf-structural-derivatives} supplies derivatives of structural
maps and parameter transforms; Chapter~\ref{ch:bf-factor-derivatives} defines
what differentiated SVD factors must reconstruct; Chapter~\ref{ch:bf-custom-gradient-wrappers}
requires the SVD score to differentiate the same implemented scalar value; and
Chapter~\ref{ch:bf-derivative-validation} supplies the finite-difference,
backend-parity, branch, and HMC-gating tests.

Let a Gaussian input at some predict or update substep be
\[
  X\sim\calN(m(\theta),P(\theta)).
\]
Let an SVD/eigendecomposition factor be written
\[
  P=U\Lambda U^\top,\qquad
  C=U\Lambda_+^{1/2},
\]
where $\Lambda_+$ denotes the eigenvalue policy actually used for point
generation: active positive eigenvalues, a smooth floor, or a hard floor.  A
fixed sigma-point rule can be written as
\begin{equation}
\label{eq:bf-svd-sp-cloud}
  \chi^{(r)} = m + C\xi^{(r)},\qquad r=0,\ldots,N-1,
\end{equation}
where the standardized offsets $\xi^{(r)}$ and weights
$w_r^{(m)},w_r^{(c)}$ are fixed by the UKF, CKF, or another deterministic
quadrature rule.  For the UKF, for example, the nonzero offsets are
$\pm\gamma e_a$.  If tuning parameters such as $\alpha$ or $\kappa$ are
estimated, the same formulas hold with additional $\dot\xi^{(r)}$,
$\ddot\xi^{(r)}$, $\dot w_r$, and $\ddot w_r$ terms.  BayesFilter's default
derivative contract treats the rule as fixed.

With fixed quadrature offsets, the point derivatives are
\begin{align}
\label{eq:bf-svd-sp-point-first}
  \dot\chi^{(r,i)} &= \dot m^{(i)}+\dot C^{(i)}\xi^{(r)},\\
\label{eq:bf-svd-sp-point-second}
  \ddot\chi^{(r,ij)} &= \ddot m^{(ij)}+\ddot C^{(ij)}\xi^{(r)}.
\end{align}
Thus all spectral difficulty enters through $\dot C$ and $\ddot C$; the
remaining recursion is ordinary chain rule and product rule.

\subsection{Map and Moment Derivatives}
\label{sec:bf-svd-sp-map-moment-derivatives}

Let $g_\theta$ denote either the transition map in the predict step or the
observation map in the update step, and define
\[
  \eta^{(r)} = g_\theta(\chi^{(r)}).
\]
Using $D_xg_\theta$ for the Jacobian in the state argument, and
$D^2_{xx}g_\theta[a,b]$ for the second directional derivative in directions
$a,b$, the propagated point derivatives are
\begin{align}
\label{eq:bf-svd-sp-map-first}
  \dot\eta^{(r,i)}
  &=
  D_xg_\theta(\chi^{(r)})\dot\chi^{(r,i)}
  +\partial_i g_\theta(\chi^{(r)}),\\
\label{eq:bf-svd-sp-map-second}
  \ddot\eta^{(r,ij)}
  &=
  D_xg_\theta(\chi^{(r)})\ddot\chi^{(r,ij)}
  +D^2_{xx}g_\theta(\chi^{(r)})
      [\dot\chi^{(r,i)},\dot\chi^{(r,j)}] \nonumber\\
  &\quad
  +D^2_{x\theta_i}g_\theta(\chi^{(r)})\dot\chi^{(r,j)}
  +D^2_{x\theta_j}g_\theta(\chi^{(r)})\dot\chi^{(r,i)}
  +\partial^2_{ij}g_\theta(\chi^{(r)}).
\end{align}
These equations are the exact place where the structural model derivatives
enter the SVD sigma-point score and Hessian.

The sigma-point mean and covariance of the transformed cloud are
\begin{align}
\label{eq:bf-svd-sp-transformed-mean}
  \bar\eta &= \sum_{r=0}^{N-1} w_r^{(m)}\eta^{(r)},\\
\label{eq:bf-svd-sp-transformed-cov}
  M_\eta &= \sum_{r=0}^{N-1}w_r^{(c)}
    \Delta_\eta^{(r)}\Delta_\eta^{(r)\top}+\Omega_\theta,
  \qquad
  \Delta_\eta^{(r)}=\eta^{(r)}-\bar\eta,
\end{align}
where $\Omega_\theta$ is the process-noise covariance in a predict step or the
measurement-noise covariance in an update step.  Their derivatives are
\begin{align}
\label{eq:bf-svd-sp-mean-first}
  \dot{\bar\eta}^{(i)}
  &=\sum_r w_r^{(m)}\dot\eta^{(r,i)},\\
\label{eq:bf-svd-sp-mean-second}
  \ddot{\bar\eta}^{(ij)}
  &=\sum_r w_r^{(m)}\ddot\eta^{(r,ij)},
\end{align}
and, with
$\dot\Delta_\eta^{(r,i)}=\dot\eta^{(r,i)}-\dot{\bar\eta}^{(i)}$,
\begin{align}
\label{eq:bf-svd-sp-cov-first}
  \dot M_\eta^{(i)}
  &=
  \sum_r w_r^{(c)}
  \left[
    \dot\Delta_\eta^{(r,i)}\Delta_\eta^{(r)\top}
    +\Delta_\eta^{(r)}\dot\Delta_\eta^{(r,i)\top}
  \right]
  +\dot\Omega_\theta^{(i)},\\
\label{eq:bf-svd-sp-cov-second}
  \ddot M_\eta^{(ij)}
  &=
  \sum_r w_r^{(c)}
  \left[
    \ddot\Delta_\eta^{(r,ij)}\Delta_\eta^{(r)\top}
    +\Delta_\eta^{(r)}\ddot\Delta_\eta^{(r,ij)\top}
    +\dot\Delta_\eta^{(r,i)}\dot\Delta_\eta^{(r,j)\top}
    +\dot\Delta_\eta^{(r,j)}\dot\Delta_\eta^{(r,i)\top}
  \right]
  +\ddot\Omega_\theta^{(ij)} .
\end{align}
Equations
\eqref{eq:bf-svd-sp-point-first}--\eqref{eq:bf-svd-sp-cov-second}
are the analytic derivative recursion for any sigma-point predict moment.

\subsection{UKF Update Objects}
\label{sec:bf-svd-sp-update-objects}

At the observation update, take the input mean and covariance to be
$m_t^-=m_{t|t-1}$ and $P_t^-=P_{t|t-1}$, construct points
$\chi_t^{(r)}$ by \eqref{eq:bf-svd-sp-cloud}, and propagate
\[
  z_t^{(r)}=h_\theta(\chi_t^{(r)}).
\]
Equations \eqref{eq:bf-svd-sp-mean-first}--\eqref{eq:bf-svd-sp-cov-second}
give
\[
  \bar z_t,\quad
  S_t=\sum_r w_r^{(c)}
    (z_t^{(r)}-\bar z_t)(z_t^{(r)}-\bar z_t)^\top+R_\theta,
\]
and their first and second derivatives.  The innovation is
\begin{align}
\label{eq:bf-svd-sp-innovation-derivatives}
  v_t &= y_t-\bar z_t, &
  \dot v_t^{(i)} &= -\dot{\bar z}_t^{(i)}, &
  \ddot v_t^{(ij)} &= -\ddot{\bar z}_t^{(ij)},
\end{align}
unless the data transformation itself depends on parameters.

The cross-covariance used in the Kalman gain is
\begin{equation}
\label{eq:bf-svd-sp-cross-cov}
  P_{xz,t}
  =
  \sum_r w_r^{(c)}
  \Delta_x^{(r)}\Delta_z^{(r)\top},
  \qquad
  \Delta_x^{(r)}=\chi_t^{(r)}-m_t^-,
  \qquad
  \Delta_z^{(r)}=z_t^{(r)}-\bar z_t .
\end{equation}
Its first and second derivatives are
\begin{align}
\label{eq:bf-svd-sp-cross-first}
  \dot P_{xz,t}^{(i)}
  &=
  \sum_r w_r^{(c)}
  \left[
    \dot\Delta_x^{(r,i)}\Delta_z^{(r)\top}
    +\Delta_x^{(r)}\dot\Delta_z^{(r,i)\top}
  \right],\\
\label{eq:bf-svd-sp-cross-second}
  \ddot P_{xz,t}^{(ij)}
  &=
  \sum_r w_r^{(c)}
  \left[
    \ddot\Delta_x^{(r,ij)}\Delta_z^{(r)\top}
    +\Delta_x^{(r)}\ddot\Delta_z^{(r,ij)\top}
    +\dot\Delta_x^{(r,i)}\dot\Delta_z^{(r,j)\top}
    +\dot\Delta_x^{(r,j)}\dot\Delta_z^{(r,i)\top}
  \right].
\end{align}
The gain is best differentiated as a solved matrix equation.  Since
$K_t=P_{xz,t}S_t^{-1}$, equivalently
\[
  S_tK_t^\top=P_{xz,t}^\top.
\]
Therefore
\begin{align}
\label{eq:bf-svd-sp-gain-first}
  \dot K_t^{(i)\top}
  &=
  \mathcal{S}_{S_t}
  \left(\dot P_{xz,t}^{(i)\top}-\dot S_t^{(i)}K_t^\top\right),\\
\label{eq:bf-svd-sp-gain-second}
  \ddot K_t^{(ij)\top}
  &=
  \mathcal{S}_{S_t}
  \left(
    \ddot P_{xz,t}^{(ij)\top}
    -\ddot S_t^{(ij)}K_t^\top
    -\dot S_t^{(i)}\dot K_t^{(j)\top}
    -\dot S_t^{(j)}\dot K_t^{(i)\top}
  \right).
\end{align}
Finally, the usual sigma-point update
\begin{align}
\label{eq:bf-svd-sp-filtered-update}
  m_t^+ &= m_t^-+K_tv_t,\\
  P_t^+ &= P_t^- - K_tS_tK_t^\top
\end{align}
has derivatives
\begin{align}
\label{eq:bf-svd-sp-filtered-mean-first}
  \dot m_t^{+(i)}
  &=
  \dot m_t^{-(i)}
  +\dot K_t^{(i)}v_t
  +K_t\dot v_t^{(i)},\\
\label{eq:bf-svd-sp-filtered-mean-second}
  \ddot m_t^{+(ij)}
  &=
  \ddot m_t^{-(ij)}
  +\ddot K_t^{(ij)}v_t
  +\dot K_t^{(i)}\dot v_t^{(j)}
  +\dot K_t^{(j)}\dot v_t^{(i)}
  +K_t\ddot v_t^{(ij)}.
\end{align}
For the covariance, define
\begin{align}
\label{eq:bf-svd-sp-filtered-cov-first}
  \dot P_t^{+(i)}
  &=
  \dot P_t^{-(i)}
  -\dot K_t^{(i)}S_tK_t^\top
  -K_t\dot S_t^{(i)}K_t^\top
  -K_tS_t\dot K_t^{(i)\top},
\end{align}
and
\begin{align}
\label{eq:bf-svd-sp-filtered-cov-second}
  \ddot P_t^{+(ij)}
  &=
  \ddot P_t^{-(ij)}
  -\ddot K_t^{(ij)}S_tK_t^\top
  -\dot K_t^{(i)}\dot S_t^{(j)}K_t^\top
  -\dot K_t^{(i)}S_t\dot K_t^{(j)\top} \nonumber\\
  &\quad
  -\dot K_t^{(j)}\dot S_t^{(i)}K_t^\top
  -K_t\ddot S_t^{(ij)}K_t^\top
  -K_t\dot S_t^{(i)}\dot K_t^{(j)\top} \nonumber\\
  &\quad
  -\dot K_t^{(j)}S_t\dot K_t^{(i)\top}
  -K_t\dot S_t^{(j)}\dot K_t^{(i)\top}
  -K_tS_t\ddot K_t^{(ij)\top}.
\end{align}
The next time step then refactors $P_t^+$ and differentiates the new SVD factor.
This is where the spectral contract below enters the recursion again.

\subsection{Likelihood Score and Hessian}
\label{sec:bf-svd-sp-likelihood-derivatives}

The SVD sigma-point likelihood contribution is the Gaussian innovation
likelihood
\[
  \ell_t
  =
  -\frac{1}{2}
  \left[
    \log\det S_t+v_t^\top S_t^{-1}v_t+n_y\log(2\pi)
  \right],
\]
where $S_t$ and $v_t$ are the sigma-point objects above.  Let
$S_tw_t=v_t$.  The score is
\begin{equation}
\label{eq:bf-svd-sp-score}
  \partial_i\ell_t
  =
  -\frac{1}{2}
  \left[
    \tr(S_t^{-1}\dot S_t^{(i)})
    +2\dot v_t^{(i)\top}w_t
    -w_t^\top\dot S_t^{(i)}w_t
  \right],
\end{equation}
and the Hessian is obtained by substituting the sigma-point
$\dot v_t,\dot S_t,\ddot v_t,\ddot S_t$ into
\eqref{eq:bf-solve-hessian-w-dot}--\eqref{eq:bf-solve-hessian-quadratic}.
Thus the UKF/SVD Hessian is not a separate likelihood theorem.  It is the
Kalman solve-form Hessian applied to approximate moments generated by
\eqref{eq:bf-svd-sp-cloud}--\eqref{eq:bf-svd-sp-filtered-cov-second}.

\section{Structural Fixed-Support Score}
\label{sec:bf-svd-sp-structural-fixed-support-score}

Chapter~\ref{ch:bf-structural-deterministic-dynamics} explains why a
structural sigma-point filter places points on the pre-transition uncertainty
variable
\begin{equation}
\label{eq:bf-svd-sp-structural-input-law}
  A_t=(x_{t-1},\varepsilon_t),
  \qquad
  A_t\mid y_{1:t-1}\approx\calN(\mu_{a,t},P_{a,t}),
\end{equation}
and then completes the state pointwise:
\begin{equation}
\label{eq:bf-svd-sp-structural-point-map}
  A_t^{(r)}=\mu_{a,t}+C_{a,t}\xi^{(r)},\qquad
  X_t^{(r)}=F_\theta(A_t^{(r)}),\qquad
  Z_t^{(r)}=h_\theta(X_t^{(r)}).
\end{equation}
Equations~\eqref{eq:bf-svd-sp-structural-input-law}--%
\eqref{eq:bf-svd-sp-structural-point-map} are the derivative version of
Proposition~\ref{prop:bf-structural-ukf-pushforward}.  The filtered state is
still $x_t$; the sigma-point variable is the pre-transition uncertainty whose
pushforward generates the predictive law.  This is the object that must be
differentiated for deterministic-completion examples such as Model C in
Chapter~\ref{ch:bf-nonlinear-ssm-validation}.

The subtle case is a fixed structural zero direction.  Write
\[
  P_{a,t}=U_t\Lambda_tU_t^\top,\qquad
  \Lambda_t=\operatorname{diag}(\lambda_{1t},\ldots,\lambda_{dt}),
\]
and split the spectral indices into
\[
  \mathcal I_t^+=\{a:\lambda_{at}>\tau\},\qquad
  \mathcal I_t^0=\{a:\lambda_{at}\le \tau\}.
\]
The structural fixed-support branch assumes that the zero-support block is
part of the model law, not a numerical floor:
\begin{align}
\label{eq:bf-svd-sp-fixed-support-null-conditions}
  u_b^\top \dot P_{a,t}^{(i)}u_c &=0
  \qquad
  \text{whenever } b\in\mathcal I_t^0
  \text{ or } c\in\mathcal I_t^0,\\
\label{eq:bf-svd-sp-fixed-support-gap-condition}
  |\lambda_{at}-\lambda_{bt}| &> g_{\min}
  \qquad
  \text{for every differentiated active pair.}
\end{align}
Here ``differentiated active pair'' means any pair that enters the derivative
of an active factor column: active-active pairs and active-null pairs.  Gaps
between two null directions do not matter because their factor columns are
identically zero on this branch.

The implemented factor is
\begin{equation}
\label{eq:bf-svd-sp-fixed-support-factor}
  C_{a,t}
  =
  \begin{bmatrix}
    \sqrt{\lambda_{1t}}u_{1t} & \cdots & \sqrt{\lambda_{dt}}u_{dt}
  \end{bmatrix},
  \qquad
  \sqrt{\lambda_{at}}u_{at}=0
  \quad\text{for }a\in\mathcal I_t^0 .
\end{equation}
For an active column $a\in\mathcal I_t^+$,
\begin{align}
\label{eq:bf-svd-sp-fixed-support-eigenvector-first}
  \dot u_{at}^{(i)}
  &=
  \sum_{b\ne a}
  u_{bt}
  \frac{u_{bt}^\top\dot P_{a,t}^{(i)}u_{at}}
       {\lambda_{at}-\lambda_{bt}},\\
\label{eq:bf-svd-sp-fixed-support-column-first}
  \dot c_{at}^{(i)}
  &=
  \sqrt{\lambda_{at}}\,\dot u_{at}^{(i)}
  +
  \frac{u_{at}^\top\dot P_{a,t}^{(i)}u_{at}}
       {2\sqrt{\lambda_{at}}}\,u_{at}.
\end{align}
For a null column $a\in\mathcal I_t^0$,
\begin{equation}
\label{eq:bf-svd-sp-fixed-support-null-column}
  c_{at}=0,\qquad \dot c_{at}^{(i)}=0.
\end{equation}
Under \eqref{eq:bf-svd-sp-fixed-support-null-conditions}, these column rules
reconstruct the derivative of the same pre-transition covariance:
\begin{equation}
\label{eq:bf-svd-sp-fixed-support-reconstruction}
  \dot P_{a,t}^{(i)}
  =
  \dot C_{a,t}^{(i)}C_{a,t}^\top
  +C_{a,t}\dot C_{a,t}^{(i)\top}.
\end{equation}
The reconstruction target is not a nuggeted covariance.  A positive floor that
turns a structural zero direction into positive variance changes the law and
must be reported as a different regularized branch.

With this fixed-support factor, the score recursion is the ordinary
sigma-point chain rule.  For a scalar parameter component $\theta_i$,
\begin{align}
\label{eq:bf-svd-sp-structural-point-first}
  \dot A_t^{(r,i)}
  &=
  \dot\mu_{a,t}^{(i)}+\dot C_{a,t}^{(i)}\xi^{(r)},\\
\label{eq:bf-svd-sp-structural-state-point-first}
  \dot X_t^{(r,i)}
  &=
  D_aF_\theta(A_t^{(r)})\dot A_t^{(r,i)}
  +\partial_iF_\theta(A_t^{(r)}),\\
\label{eq:bf-svd-sp-structural-observation-point-first}
  \dot Z_t^{(r,i)}
  &=
  D_xh_\theta(X_t^{(r)})\dot X_t^{(r,i)}
  +\partial_i h_\theta(X_t^{(r)}).
\end{align}
The moment derivatives are
\begin{align}
\label{eq:bf-svd-sp-structural-moment-first}
  \dot{\hat x}_t^{(i)}
    &=\sum_r w_r^{(m)}\dot X_t^{(r,i)},&
  \dot{\hat z}_t^{(i)}
    &=\sum_r w_r^{(m)}\dot Z_t^{(r,i)},\\
\label{eq:bf-svd-sp-structural-S-first}
  \dot S_t^{(i)}
    &=
    \sum_r w_r^{(c)}
    \left[
      \dot\Delta_z^{(r,i)}\Delta_z^{(r)\top}
      +
      \Delta_z^{(r)}\dot\Delta_z^{(r,i)\top}
    \right]
    +\dot R_t^{(i)},\\
\label{eq:bf-svd-sp-structural-cross-first}
  \dot P_{xz,t}^{(i)}
    &=
    \sum_r w_r^{(c)}
    \left[
      \dot\Delta_x^{(r,i)}\Delta_z^{(r)\top}
      +
      \Delta_x^{(r)}\dot\Delta_z^{(r,i)\top}
    \right],
\end{align}
where
\[
  \Delta_x^{(r)}=X_t^{(r)}-\hat x_t,\qquad
  \Delta_z^{(r)}=Z_t^{(r)}-\hat z_t,
\]
and dotted centered quantities subtract the corresponding dotted means.  The
innovation derivative remains
\[
  v_t=y_t-\hat z_t,\qquad
  \dot v_t^{(i)}=-\dot{\hat z}_t^{(i)}.
\]
Therefore the per-period score is
\begin{equation}
\label{eq:bf-svd-sp-structural-score}
  \partial_i\ell_t
  =
  -\frac12
  \left[
    \operatorname{tr}(S_t^{-1}\dot S_t^{(i)})
    +2\dot v_t^{(i)\top}S_t^{-1}v_t
    -v_t^\top S_t^{-1}\dot S_t^{(i)}S_t^{-1}v_t
  \right].
\end{equation}
This sign convention is the direct derivative of the likelihood in
\eqref{eq:bf-structural-ukf-loglik-object}; the term
$\dot v_t=-\dot{\hat z}_t$ is where the observation mean derivative enters.

The diagnostics for this branch must report at least the fixed null count, the
maximum fixed-null derivative residual in
\eqref{eq:bf-svd-sp-fixed-support-null-conditions}, the active spectral gap in
\eqref{eq:bf-svd-sp-fixed-support-gap-condition}, the factor reconstruction
residual in \eqref{eq:bf-svd-sp-fixed-support-reconstruction}, and the
deterministic-completion residual of the propagated points.  Passing these
checks means that the score differentiates the implemented structural
sigma-point likelihood.  It does not certify a nonlinear Hessian, HMC target,
or GPU/XLA scaling claim.

\section{Spectral Factor Derivative Contract}
\label{sec:bf-svd-factor-derivative-contract}

The preceding section assumes $\dot C$ and $\ddot C$ are available.  This
subsection states what those derivatives mean for an SVD/eigen factor.  Suppose
first that $P(\theta)$ has a smooth simple spectrum and no active hard floor:
\[
  P=\sum_{a=1}^n \lambda_a u_au_a^\top,\qquad
  \lambda_a>0,\qquad \lambda_a\ne\lambda_b\ (a\ne b).
\]
For a parameter component $\theta_i$,
\begin{align}
\label{eq:bf-svd-eigenvalue-first}
  \dot\lambda_a^{(i)}
  &=u_a^\top\dot P^{(i)}u_a,\\
\label{eq:bf-svd-eigenvector-first}
  \dot u_a^{(i)}
  &=
  \sum_{b\ne a}
  u_b\,
  \frac{u_b^\top\dot P^{(i)}u_a}{\lambda_a-\lambda_b}.
\end{align}
For the column factor $c_a=\sqrt{\lambda_a}\,u_a$,
\begin{equation}
\label{eq:bf-svd-factor-column-first}
  \dot c_a^{(i)}
  =
  \sqrt{\lambda_a}\,\dot u_a^{(i)}
  +\frac{\dot\lambda_a^{(i)}}{2\sqrt{\lambda_a}}u_a .
\end{equation}
Second derivatives can be written in the same eigenbasis.  Define
$A_a^{(i)}=\dot P^{(i)}-\dot\lambda_a^{(i)}I$ and
$A_a^{(ij)}=\ddot P^{(ij)}-\ddot\lambda_a^{(ij)}I$.  Then
\begin{equation}
\label{eq:bf-svd-eigenvalue-second}
  \ddot\lambda_a^{(ij)}
  =
  u_a^\top\ddot P^{(ij)}u_a
  +2\sum_{b\ne a}
    \frac{
      (u_b^\top\dot P^{(i)}u_a)
      (u_b^\top\dot P^{(j)}u_a)
    }{\lambda_a-\lambda_b}.
\end{equation}
The second derivative of the eigenvector is determined by differentiating the
eigen-equation and the normalization condition:
\begin{equation}
\label{eq:bf-svd-eigenvector-second}
  \ddot u_a^{(ij)}
  =
  -u_a\,\dot u_a^{(i)\top}\dot u_a^{(j)}
  +\sum_{b\ne a}u_b
  \frac{
    u_b^\top
    \left(
      A_a^{(ij)}u_a
      +A_a^{(i)}\dot u_a^{(j)}
      +A_a^{(j)}\dot u_a^{(i)}
    \right)
  }{\lambda_a-\lambda_b}.
\end{equation}
Consequently
\begin{align}
\label{eq:bf-svd-factor-column-second}
  \ddot c_a^{(ij)}
  &=
  \sqrt{\lambda_a}\,\ddot u_a^{(ij)}
  +\frac{\dot\lambda_a^{(i)}}{2\sqrt{\lambda_a}}\dot u_a^{(j)}
  +\frac{\dot\lambda_a^{(j)}}{2\sqrt{\lambda_a}}\dot u_a^{(i)}
  \nonumber\\
  &\quad
  +\left(
      \frac{\ddot\lambda_a^{(ij)}}{2\sqrt{\lambda_a}}
      -\frac{\dot\lambda_a^{(i)}\dot\lambda_a^{(j)}}{4\lambda_a^{3/2}}
    \right)u_a .
\end{align}
Equations
\eqref{eq:bf-svd-eigenvalue-first}--\eqref{eq:bf-svd-factor-column-second}
are the branch-local eigenderivative formulas for an eigen-factor covariance
square root on a smooth simple-spectrum branch, before any active hard-floor
or ordering/sign fallback is crossed.

The formulas also show why the derivative path is fragile.  Eigenvector
derivatives contain denominators $\lambda_a-\lambda_b$; the Hessian contains
the same gaps again through $\dot u_a$ and through differentiation of the gap.
This is the spectral-gap risk identified by matrix-backpropagation calculus
\citep{ionescu2015matrixbackprop}.  If eigenvalues are repeated, individual
eigenvectors are not uniquely differentiable.  If a hard floor
$\lambda_a^+=\max(\lambda_a,\epsilon)$ is used, the derivative is branch
dependent and undefined at $\lambda_a=\epsilon$.  A smooth floor
$\lambda_a^+=\phi_\epsilon(\lambda_a)$ replaces
$\dot\lambda_a$ and $\ddot\lambda_a$ by
\[
  \dot\lambda_a^+=\phi_\epsilon'(\lambda_a)\dot\lambda_a,\qquad
  \ddot\lambda_a^+
  =
  \phi_\epsilon'(\lambda_a)\ddot\lambda_a
  +\phi_\epsilon''(\lambda_a)\dot\lambda_a^{(i)}\dot\lambda_a^{(j)},
\]
but it does not remove the eigenvector gap denominators.

For this reason BayesFilter should validate SVD factor derivatives through
reconstruction identities, not merely by checking that autodiff returns finite
numbers.  The reconstruction target is the covariance represented by the
implemented factor,
\[
  P_+ = CC^\top=U\Lambda_+U^\top,
\]
not necessarily the pre-floor covariance $P$:
\begin{align}
\label{eq:bf-svd-factor-reconstruction-first}
  \dot P_+^{(i)}
  &\stackrel{?}{=}
  \dot C^{(i)}C^\top+C\dot C^{(i)\top},\\
\label{eq:bf-svd-factor-reconstruction-second}
  \ddot P_+^{(ij)}
  &\stackrel{?}{=}
  \ddot C^{(ij)}C^\top+C\ddot C^{(ij)\top}
  +\dot C^{(i)}\dot C^{(j)\top}
  +\dot C^{(j)}\dot C^{(i)\top}.
\end{align}
The residuals in
\eqref{eq:bf-svd-factor-reconstruction-first}--\eqref{eq:bf-svd-factor-reconstruction-second},
the minimum spectral gap, the active-floor set, and any fallback branch are part
of the derivative output contract.

\section{Principal-Square-Root Placement Backend}
\label{sec:bf-svd-principal-sqrt-backend}

The preceding eigenderivative formulas define the historical
\texttt{tf\_svd\_ukf} branch.  BayesFilter also has a distinct experimental
backend, \texttt{tf\_principal\_sqrt\_ukf}, whose factor is the principal
matrix square root rather than a signed or ordered eigenvector column factor.
The two names are backend contracts, not aliases.

On a strict-SPD branch, let
\[
  P\in\mathbb S_{++}^n,\qquad
  S=P^{1/2}=S^\top\succ0,\qquad
  S^2=P .
\]
For a covariance perturbation $dP$, the Frechet derivative of the principal
square root is the unique solution $dS$ of the Sylvester equation
\begin{equation}
\label{eq:bf-svd-principal-sqrt-sylvester-first}
  S\,dS+dS\,S=dP .
\end{equation}
Equivalently, for a parameter component $\theta_i$,
\begin{equation}
\label{eq:bf-svd-principal-sqrt-first}
  S\,\dot S^{(i)}+\dot S^{(i)}S=\dot P^{(i)} .
\end{equation}
Because the eigenvalues of strict-SPD $S$ are positive, the linear
Sylvester operator $X\mapsto SX+XS$ is invertible on all matrices.  This is the
reason the principal-square-root derivative does not contain the active
repeated-positive-eigenvalue denominators that appear in
\eqref{eq:bf-svd-eigenvector-first}.  This statement is only a strict-SPD
square-root derivative statement.  It is not a rank-loss theorem, not a
structural-null completion rule, and not an HMC convergence result.

For second derivatives on the same strict-SPD branch, differentiate
\eqref{eq:bf-svd-principal-sqrt-first} once more:
\begin{equation}
\label{eq:bf-svd-principal-sqrt-second}
  S\,\ddot S^{(ij)}+\ddot S^{(ij)}S
  =
  \ddot P^{(ij)}
  -\dot S^{(i)}\dot S^{(j)}
  -\dot S^{(j)}\dot S^{(i)} .
\end{equation}
Thus the same Sylvester primitive supplies the Hessian-level square-root
factor derivative once $\dot S^{(i)}$, $\dot S^{(j)}$, and
$\ddot P^{(ij)}$ are available.

The sigma-point rule has square-root gauge freedom.  If $C C^\top=P$, then
the standardized point cloud $\chi^{(r)}=m+C\xi^{(r)}$ matches the same first
and second Gaussian moments whenever the chosen rule's offsets and weights
have the required standardized moment identities.  Replacing an eigenvector
factor $C=U\Lambda^{1/2}$ by the principal factor $S=P^{1/2}$ is therefore a
change of square-root gauge for moment placement.  It does not imply that
nonlinear propagated values may differ only by roundoff: after a nonlinear map
$g$, the clouds $\{g(m+C\xi^{(r)})\}$ and $\{g(m+S\xi^{(r)})\}$ may differ.
BayesFilter must therefore validate the implemented likelihood and score, not
claim nonlinear value identity across gauges.

The current principal-square-root backend is strict-SPD only.  It must keep
fixed-null support disabled unless a separate reviewed branch specifies how
structural nulls, active floors, and rank-deficient covariance laws are
represented and differentiated.  Required diagnostics for this backend include
the Sylvester residual in
\eqref{eq:bf-svd-principal-sqrt-first}, the reconstruction residual
$\dot P^{(i)}-(\dot S^{(i)}S+S\dot S^{(i)})$, minimum eigenvalue of $P$, any
jitter or floor policy, and eager/compiled parity for the complete
value-and-score path.

The Phase 2026-06 strict-SPD implementation evidence is bounded: repeated
positive-spectrum fixtures, reconstruction checks, compiled value/score smoke,
and the two-candidate NK replay are engineering artifacts.  They do not
establish posterior correctness, HMC convergence, runtime superiority,
GPU readiness, default readiness, or full baseline launch readiness.

\section{Validation Targets for SVD Score and Hessian}
\label{sec:bf-svd-sp-derivative-validation}

The analytical derivation above leads to concrete tests:
\begin{enumerate}[label=(\roman*)]
  \item On affine Gaussian systems, the SVD sigma-point score and Hessian must
    match the exact Kalman score and Hessian from
    Chapters~\ref{ch:bf-kalman-score}--\ref{ch:bf-kalman-hessian}.
  \item On small nonlinear systems, the analytical score and Hessian must match
    finite differences of the same SVD sigma-point likelihood.
  \item The factor residuals
    \eqref{eq:bf-svd-factor-reconstruction-first}--\eqref{eq:bf-svd-factor-reconstruction-second}
    must remain small on well-separated spectra.
  \item Tests must deliberately include clustered eigenvalues, zero eigenvalues,
    floor crossings, and sign/order changes, and report when the derivative is a
    branch derivative rather than the derivative of a smooth mathematical SVD.
  \item HMC diagnostics must record whether gradients came from the structural
    analytic recursion, a custom spectral derivative, raw autodiff, or a
    regularized fallback.
\end{enumerate}
These tests turn the phrase ``SVD sigma-point Hessian'' into an auditable object:
a sigma-point moment derivative recursion, a solve-form likelihood Hessian, and
a declared spectral-factor derivative policy.

\section{SVD-CUT Score and Hessian}
\label{sec:bf-cut-svd-score-hessian}

The conjugate unscented transform can be inserted into the same spectral
sigma-point filter.  The precise object is an SVD-placed CUT rule: use an
SVD/eigen factor to place the CUT offsets, then use the ordinary sigma-point
moment and likelihood recursions.  This is sometimes called CUT-SVD and
sometimes SVD-CUT.  The name is less important than the contract
\[
  P_\star=C C^\top ,
\]
where $P_\star$ is the covariance actually represented by the implemented
factor.  If eigenvalues have been floored, truncated, or otherwise regularized,
$P_\star$ is the floored or truncated covariance, not necessarily the
pre-regularized covariance submitted to the factor routine.

Let the SVD branch return a factor $C(\theta)\in\mathbb R^{n_x\times q}$ and
let $Z\sim\calN(0,I_q)$.  A CUT rule in the factor coordinates supplies
standardized offsets and weights
\[
  \{(\xi^{(r)},w_r): r=0,\ldots,N_{\mathrm{CUT}}-1\},
  \qquad
  \xi^{(r)}\in\mathbb R^q.
\]
For CUT4-G, when $q\ge3$,
\[
  N_{\mathrm{CUT4}}=2q+2^q .
\]
The SVD-CUT points are
\begin{equation}
\label{eq:bf-svd-cut-cloud}
  \chi^{(r)}=m+C\xi^{(r)},\qquad
  r=0,\ldots,N_{\mathrm{CUT}}-1.
\end{equation}
The quadrature approximation to an expectation under the implemented Gaussian
law $X_\star=m+CZ\sim\calN(m,P_\star)$ is
\begin{equation}
\label{eq:bf-svd-cut-quadrature}
  \mathcal Q_{\mathrm{SVD}\text{-}\mathrm{CUT}}[\varphi]
  =
  \sum_{r=0}^{N_{\mathrm{CUT}}-1} w_r
  \varphi(m+C\xi^{(r)}).
\end{equation}

\begin{proposition}[Polynomial exactness under affine SVD placement]
\label{prop:bf-svd-cut-affine-exactness}
Suppose the standardized CUT rule is exact for standard-normal polynomials on
$\mathbb R^q$ through degree $d$.  If $C C^\top=P_\star$, then
\eqref{eq:bf-svd-cut-quadrature} integrates every polynomial
$p:\mathbb R^{n_x}\to\mathbb R$ of total degree at most $d$ exactly under
$X_\star\sim\calN(m,P_\star)$.  The statement remains true when
$P_\star$ is singular.
\end{proposition}

\begin{proof}
For any polynomial $p$ of degree at most $d$, define
\[
  \psi(z)=p(m+Cz).
\]
The affine composition cannot increase total degree, so $\psi$ is a polynomial
on $\mathbb R^q$ of degree at most $d$.  By standardized CUT exactness,
\[
  \sum_r w_r \psi(\xi^{(r)})
  =
  \mathbb E[\psi(Z)].
\]
Substituting the definition of $\psi$ gives
\[
  \sum_r w_r p(m+C\xi^{(r)})
  =
  \mathbb E[p(m+CZ)]
  =
  \mathbb E[p(X_\star)].
\]
No inverse of $C$ is used.  Hence rank deficiency of $P_\star=C C^\top$ does
not affect the formal quadrature exactness statement; it only changes the
affine support on which the degenerate Gaussian law lives.
\end{proof}

This proposition is the clean mathematical answer to the accuracy question:
abstracting from finite arithmetic, SVD placement does not by itself destroy
CUT accuracy.  It is not a recommendation to hide a structural state split
inside a large singular covariance.  A full zero-padded factor may reproduce
the same law, but it also creates duplicate or nearly duplicate points in null
directions and pushes the derivative path through the larger spectral branch.

\subsection{SVD-CUT Moment Derivatives}
\label{sec:bf-svd-cut-moment-derivatives}

The analytic gradient and Hessian follow by differentiating
\eqref{eq:bf-svd-cut-cloud}--\eqref{eq:bf-svd-cut-quadrature}.  Assume first
that the CUT offsets and weights are fixed.  If CUT tuning parameters are
estimated, the same derivation applies after adding the explicit
$\dot \xi^{(r)}$, $\ddot \xi^{(r)}$, $\dot w_r$, and $\ddot w_r$ terms.

For a parameter component $\theta_i$, the point derivatives are
\begin{align}
\label{eq:bf-svd-cut-point-first}
  \dot\chi^{(r,i)}
  &=
  \dot m^{(i)}+\dot C^{(i)}\xi^{(r)},\\
\label{eq:bf-svd-cut-point-second}
  \ddot\chi^{(r,ij)}
  &=
  \ddot m^{(ij)}+\ddot C^{(ij)}\xi^{(r)}.
\end{align}
Let $g_\theta$ be the nonlinear map being integrated.  In a predict step this
is the transition map; in an update step it is the observation map.  Write
\[
  y^{(r)}=g_\theta(\chi^{(r)}),\qquad
  \bar y=\sum_r w_r y^{(r)},\qquad
  d^{(r)}=y^{(r)}-\bar y .
\]
The propagated point derivatives are
\begin{align}
\label{eq:bf-svd-cut-map-first}
  \dot y^{(r,i)}
  &=
  D_xg_\theta(\chi^{(r)})\dot\chi^{(r,i)}
  +\partial_i g_\theta(\chi^{(r)}),\\
\label{eq:bf-svd-cut-map-second}
  \ddot y^{(r,ij)}
  &=
  D_xg_\theta(\chi^{(r)})\ddot\chi^{(r,ij)}
  +D^2_{xx}g_\theta(\chi^{(r)})
    [\dot\chi^{(r,i)},\dot\chi^{(r,j)}] \nonumber\\
  &\quad
  +D^2_{x\theta_i}g_\theta(\chi^{(r)})\dot\chi^{(r,j)}
  +D^2_{x\theta_j}g_\theta(\chi^{(r)})\dot\chi^{(r,i)}
  +\partial^2_{ij}g_\theta(\chi^{(r)}).
\end{align}
Therefore
\begin{align}
\label{eq:bf-svd-cut-mean-first}
  \dot{\bar y}^{(i)}
  &=
  \sum_r w_r\dot y^{(r,i)},\\
\label{eq:bf-svd-cut-mean-second}
  \ddot{\bar y}^{(ij)}
  &=
  \sum_r w_r\ddot y^{(r,ij)}.
\end{align}
Let $\Omega_\theta$ denote the additive covariance in the moment being formed:
$Q_\theta$ in a predict step or $R_\theta$ in an observation update.  The
SVD-CUT covariance moment is
\begin{equation}
\label{eq:bf-svd-cut-cov}
  M=\sum_r w_r d^{(r)}d^{(r)\top}+\Omega_\theta .
\end{equation}
With
\[
  \dot d^{(r,i)}=\dot y^{(r,i)}-\dot{\bar y}^{(i)},\qquad
  \ddot d^{(r,ij)}=\ddot y^{(r,ij)}-\ddot{\bar y}^{(ij)},
\]
its derivatives are
\begin{align}
\label{eq:bf-svd-cut-cov-first}
  \dot M^{(i)}
  &=
  \sum_r w_r
  \left[
    \dot d^{(r,i)}d^{(r)\top}
    +d^{(r)}\dot d^{(r,i)\top}
  \right]
  +\dot\Omega_\theta^{(i)},\\
\label{eq:bf-svd-cut-cov-second}
  \ddot M^{(ij)}
  &=
  \sum_r w_r
  \left[
    \ddot d^{(r,ij)}d^{(r)\top}
    +d^{(r)}\ddot d^{(r,ij)\top}
    +\dot d^{(r,i)}\dot d^{(r,j)\top}
    +\dot d^{(r,j)}\dot d^{(r,i)\top}
  \right]
  +\ddot\Omega_\theta^{(ij)}.
\end{align}

\begin{proof}[Derivation]
Equations \eqref{eq:bf-svd-cut-point-first} and
\eqref{eq:bf-svd-cut-point-second} are the first and second derivatives of the
affine point map $\chi^{(r)}=m+C\xi^{(r)}$, using fixed $\xi^{(r)}$.  Equations
\eqref{eq:bf-svd-cut-map-first} and \eqref{eq:bf-svd-cut-map-second} are the
first and second multivariate chain rules for
$g_\theta(\chi^{(r)}(\theta))$.  Equations
\eqref{eq:bf-svd-cut-mean-first} and \eqref{eq:bf-svd-cut-mean-second} follow
from linearity of the weighted sum.  Finally, differentiating
$d^{(r)}d^{(r)\top}$ once and twice gives
\[
  \partial_i(dd^\top)=\dot d_i d^\top+d\dot d_i^\top
\]
and
\[
  \partial_{ij}^2(dd^\top)
  =
  \ddot d_{ij}d^\top+d\ddot d_{ij}^\top
  +\dot d_i\dot d_j^\top+\dot d_j\dot d_i^\top ,
\]
which proves \eqref{eq:bf-svd-cut-cov-first} and
\eqref{eq:bf-svd-cut-cov-second} after summing over the fixed weights and
adding the derivatives of $\Omega_\theta$.
\end{proof}

\subsection{SVD-CUT Likelihood Gradient and Hessian}
\label{sec:bf-svd-cut-likelihood-derivatives}

For an observation update, set $g_\theta=h_\theta$ and identify
\[
  \bar z_t=\bar y,\qquad
  S_t=M,\qquad
  v_t=y_t^{\mathrm{obs}}-\bar z_t .
\]
If the recorded observation does not depend on $\theta$, then
\begin{equation}
\label{eq:bf-svd-cut-innovation-derivatives}
  \dot v_t^{(i)}=-\dot{\bar z}_t^{(i)},\qquad
  \ddot v_t^{(ij)}=-\ddot{\bar z}_t^{(ij)} .
\end{equation}
Let $S_tw_t=v_t$.  The per-time SVD-CUT log likelihood is
\begin{equation}
\label{eq:bf-svd-cut-loglik}
  \ell_t
  =
  -\frac{1}{2}
  \left[
    \log\det S_t+v_t^\top S_t^{-1}v_t+n_y\log(2\pi)
  \right].
\end{equation}
Its score is
\begin{equation}
\label{eq:bf-svd-cut-score}
  \partial_i\ell_t
  =
  -\frac{1}{2}
  \left[
    \tr(S_t^{-1}\dot S_t^{(i)})
    +2\dot v_t^{(i)\top}w_t
    -w_t^\top\dot S_t^{(i)}w_t
  \right].
\end{equation}
For the Hessian, define
\begin{equation}
\label{eq:bf-svd-cut-w-dot}
  \dot w_t^{(j)}
  =
  \mathcal{S}_{S_t}
  \left(\dot v_t^{(j)}-\dot S_t^{(j)}w_t\right).
\end{equation}
Then
\begin{align}
\label{eq:bf-svd-cut-hessian-pieces}
  T_{ij}^{\log\det}
  &=
  \tr\!\left(
    S_t^{-1}\ddot S_t^{(ij)}
    -S_t^{-1}\dot S_t^{(j)}S_t^{-1}\dot S_t^{(i)}
  \right),\\
  T_{ij}^{v}
  &=
  2\ddot v_t^{(ij)\top}w_t
  +2\dot v_t^{(i)\top}\dot w_t^{(j)},\\
  T_{ij}^{S}
  &=
  2\dot w_t^{(j)\top}\dot S_t^{(i)}w_t
  +w_t^\top\ddot S_t^{(ij)}w_t ,
\end{align}
and
\begin{equation}
\label{eq:bf-svd-cut-hessian}
  \partial_{ij}^2\ell_t
  =
  -\frac{1}{2}
  \left[
    T_{ij}^{\log\det}+T_{ij}^{v}-T_{ij}^{S}
  \right].
\end{equation}

\begin{proof}[Derivation]
Equations \eqref{eq:bf-svd-cut-cov-first} and
\eqref{eq:bf-svd-cut-cov-second}, with $\Omega_\theta=R_\theta$, supply
$\dot S_t$ and $\ddot S_t$.  Equation
\eqref{eq:bf-svd-cut-innovation-derivatives} supplies $\dot v_t$ and
$\ddot v_t$.  Differentiating the solve $S_tw_t=v_t$ gives
\[
  S_t\dot w_t^{(j)}+\dot S_t^{(j)}w_t=\dot v_t^{(j)},
\]
hence \eqref{eq:bf-svd-cut-w-dot}.  The differential identities
$\partial_i\log\det S=\tr(S^{-1}\dot S_i)$ and
$\partial_i S^{-1}=-S^{-1}\dot S_iS^{-1}$ give
\eqref{eq:bf-svd-cut-score}.  Differentiating the three pieces of
\eqref{eq:bf-svd-cut-loglik} a second time gives
\eqref{eq:bf-svd-cut-hessian-pieces} and
\eqref{eq:bf-svd-cut-hessian}.  This is the same solve-form likelihood
Hessian as Chapter~\ref{ch:bf-kalman-hessian}; the SVD-CUT-specific work is the
construction of $v_t,S_t,\dot v_t,\dot S_t,\ddot v_t,\ddot S_t$ from the CUT
point cloud.
\end{proof}

All dependence on the SVD backend enters through
$C,\dot C,\ddot C$.  Consequently the derivative must reconstruct the same
implemented covariance branch:
\[
  \dot P_\star^{(i)}
  =
  \dot C^{(i)}C^\top+C\dot C^{(i)\top},
\]
and
\[
  \ddot P_\star^{(ij)}
  =
  \ddot C^{(ij)}C^\top+C\ddot C^{(ij)\top}
  +\dot C^{(i)}\dot C^{(j)\top}
  +\dot C^{(j)}\dot C^{(i)\top}.
\]
These are exactly the reconstruction identities
\eqref{eq:bf-svd-factor-reconstruction-first}--\eqref{eq:bf-svd-factor-reconstruction-second}.
If a hard eigenvalue floor is active, the equations describe the branch
derivative of the floored covariance $P_\star$, not the derivative of the
pre-floor covariance.  If eigenvalues are clustered, repeated, reordered, or
sign-flipped, the smooth-branch formulas in
Section~\ref{sec:bf-svd-factor-derivative-contract} may fail or become
ill-conditioned.

\subsection{Point Count, GPU, and XLA}
\label{sec:bf-svd-cut-gpu-xla}

Modern GPUs and XLA make large sigma-point clouds much less painful than
scalar Python loops, but they do not make the point count disappear.  The CUT
map evaluations are embarrassingly parallel over points and chains, and the
point dimension is static when $q$ is fixed.  This fits the XLA discipline in
Chapter~\ref{ch:bf-xla-jit}: construct a fixed tensor of offsets, evaluate the
transition or observation map in batch, and reduce over the point axis inside
the compiled value-and-gradient path.

The remaining costs still scale with
\[
  N_{\mathrm{CUT4}}=2q+2^q .
\]
For one observation update with output dimension $n_y$, the value-side moment
assembly contains
\[
  O(N_{\mathrm{CUT}}\,c_h)
  \quad\text{map work},\qquad
  O(N_{\mathrm{CUT}}\,n_y^2)
  \quad\text{innovation-covariance reductions},
\]
plus
$O(N_{\mathrm{CUT}}\,n_x n_y)$ if the cross covariance and Kalman gain are
formed.  With $p$ parameters, a direct analytic score stores or recomputes
$O(N_{\mathrm{CUT}}\,p\,n_y)$ propagated first derivatives, and a direct
Hessian stores or recomputes
$O(N_{\mathrm{CUT}}\,p^2\,n_y)$ propagated second derivatives, in addition to
the SVD-factor derivative work.  XLA can fuse many elementwise operations and
run the point axis in parallel, but memory traffic, reductions, spectral
factorization, compile time, and the sequential time scan remain real
constraints.

Thus the correct engineering conclusion is conditional.  A large number of
CUT points may be acceptable on modern GPU/XLA hardware when the stochastic
factor dimension $q$ is modest, the maps are sufficiently expensive to amortize
compile and launch overhead, tensor shapes are static, and chains or particles
provide enough batch parallelism.  It is not generally acceptable to collapse a
large DSGE state into a full-state CUT rule and rely on the GPU to absorb the
$2^q$ corner set.  The preferred design remains structural SVD-CUT: apply
higher-order CUT quadrature only to the declared stochastic integration block,
then complete deterministic coordinates through the structural transition map
and use SVD regularization only where the covariance factor actually needs it.

\section{Status Labels for SVD-HMC}
\label{sec:bf-svd-status-labels}

The SVD/HMC gap-closure plan introduced four labels that BayesFilter should
reuse:
\begin{description}
  \item[filter-correct] the likelihood and gradients agree with a reference
    behavior on controlled cases;
  \item[sampler-usable] short or medium HMC runs produce finite chains and
    useful diagnostics;
  \item[converged] strict multi-chain posterior diagnostics pass on the
    intended target;
  \item[production-XLA-gated] fixed-solution, full-solve, boundary rejection,
    gradient, and tiny HMC paths compile and run under the production
    compilation policy.
\end{description}
These labels prevent a recurring mistake: a clean smoke run is not a
convergence result, and a finite likelihood is not a validated HMC gradient.

\section{BayesFilter Policy}
\label{sec:bf-svd-policy}

BayesFilter may use SVD sigma-point filters as robust approximate likelihood
infrastructure and as a diagnostic backend.  To use them as HMC production
targets, the project must additionally provide:
\begin{enumerate}[label=(\roman*)]
  \item affine-Gaussian recovery against the exact Kalman backend;
  \item nonlinear controlled recovery against dense quadrature or a trusted
    reference;
  \item finite value and finite gradient stress tests across invalid and
    near-boundary parameter proposals;
  \item eigenvalue or singular-value gap telemetry during HMC;
  \item eager/compiled parity for the complete value-and-gradient path;
  \item a custom-gradient or analytic-gradient policy if raw spectral autodiff
    fails the preceding checks.
\end{enumerate}
For industrial-size models, the conservative default is prevention rather than
debugging: certify a derivative path before relying on SVD-tape gradients in a
long HMC run.
